% Source: Weinberg GR, physical PDF pp. 439--449 (printed pp. 419--427),
%         together with source-supplements/printed-p418.png (printed p. 418).
% Physical pp. 444--445 repeat physical pp. 442--443 (printed pp. 422--423)
% and are transcribed only once.
% Coverage: Section 14.4, Measures of Distance; equations
%           (14.4.1)--(14.4.45), all intervening unnumbered displays, and the
%           concluding photometric conventions before Section 14.5.
% Figures/tables/footnotes: Figures 14.1--14.2; linked reference markers
%                          14--16; no tables or footnotes.
% Status: source-reviewed and compile-clean; visually proofed.
% Uncertainties: the restored leaf, repeated leaves, scale-factor, light-path
%                parameter, parallax, Boltzmann-constant, and source wording
%                modernizations are documented below.
%
% RESTORED SOURCE: source-supplements/printed-p418.png supplies the beginning
% of this section before physical p. 439; it was checked at full resolution.
% MODERNIZATION CHECK: R_source(t) -> a(t). The source path parameter rho is
% renamed sigma to avoid collision with the cosmic density rho; its small
% parallax angle theta is written vartheta to avoid collision with the polar
% coordinate; trigonometric parallax pi is written varpi; and Boltzmann's
% constant is k_B, reserving k for spatial curvature.
% VERIFY SOURCE: after (14.4.37) the scan calls l a bolometric apparent
% “magnitude,” but (14.4.12) defines it as luminosity. The prose below uses
% “luminosity” and records the printed slip here rather than silently hiding it.
% LITERAL-OPERATOR AUDIT: every sign in (14.4.2)--(14.4.8), the two red-shift
% factors in (14.4.12), and the exponents and -1 terms in (14.4.34)--(14.4.38)
% were checked directly against the source images.

\subsection{Measures of Distance}\label{sec:14.4}

There are at present only two practical methods (not counting the measurement
of red shifts) for determining the distance of an object outside our galaxy.
If we know its absolute luminosity, we can compare it to the observed apparent
luminosity; or if we know its true diameter, we can compare it to the observed
angular diameter. In addition, the distance to a near enough object can be
determined by measuring its \emph{parallax}, the shift in apparent position in
the sky caused by the earth's revolution about the sun, or its \emph{proper
motion}, the shift in apparent position in the sky caused by the object's
actual motion relative to the sun.

The ``distances'' measured by these four methods are identical for objects
that are nearer than about \(10^9\) light years, but beyond this range they
differ from each other, and also from the ``proper distance'' defined in
Section~\ref{sec:14.2}. Thus, in order to use the correlation between red
shifts and apparent luminosities or angular diameters to measure \(a(t)\) and
\(k\), it will first be necessary to express the distances determined from
apparent luminosities or angular diameters in terms of \(r_1\) and \(t_0\).
It will be instructive, if largely academic, to do the same for distances
determined from measurements of parallax or proper motion.

In order to calculate parallaxes and apparent luminosities, we must know the
paths of light rays that leave a source at \(r_1,\theta_1,\phi_1\) and pass
near \(r=0\). (See Figure~\ref{fig:14.1}.) In a coordinate system
\(x'^\mu\) in which the light source is at the origin, the ray path is given
by the very simple equation
\begin{equation}
  \mathbf{x}'(\sigma)=\mathbf{n}\sigma,
  \tag{14.4.1}\label{eq:14.4.1}
\end{equation}
where \(\mathbf{n}\) is a fixed vector, \(\sigma\) is a variable positive
parameter describing positions along the path (with \(\sigma=0\) at the
source), and \(\mathbf{x}'\) is a three-vector formed from the comoving
coordinates \(r',\theta',\phi'\) in the usual way:
\[
  \mathbf{x}'
  \equiv
  \left(
    r'\sin\theta'\cos\phi',
    r'\sin\theta'\sin\phi',
    r'\cos\theta'
  \right).
\]
The transformation between the \(x'^\mu\) coordinates and another coordinate
system in which the light source is at \(\mathbf{x}_1\) is given by setting
\(\mathbf{b}=\mathbf{x}_1\) and interchanging \(\mathbf{x}\) and
\(\mathbf{x}'\) in Eq.~\eqref{eq:14.2.7}:
\begin{equation}
  \mathbf{x}
  =
  \mathbf{x}'
  +\mathbf{x}_1
  \left[
    \sqrt{1-k\mathbf{x}'^2}
    -
    \left(1-\sqrt{1-k\mathbf{x}_1^2}\right)
    \frac{\mathbf{x}'\cdot\mathbf{x}_1}{\mathbf{x}_1^2}
  \right].
  \tag{14.4.2}\label{eq:14.4.2}
\end{equation}

% Source: Weinberg GR, physical PDF p. 439 (printed p. 419).
% Coverage: Figure 14.1, light-ray geometry for parallax and luminosity.
% Status: vector reconstruction; source-reviewed and compile-clean.
% Uncertainties: source caption states that angles and curvature are
%                deliberately exaggerated.

\begingroup
\renewcommand{\thefigure}{14.1}
\begin{figure}[htbp]
\centering
\begin{tikzpicture}[x=1cm,y=1cm,>=Latex]
  \coordinate (O) at (0,0);
  \coordinate (S) at (0,6.4);
  \coordinate (B) at (1.15,0.45);
  \coordinate (E) at (1.75,6.0);

  \draw[line width=0.55pt] (O)--(S);
  \draw[line width=0.8pt]
    (E) .. controls (1.25,4.2) and (0.55,1.7) .. (B);
  \draw[dashed] (E)--(S);
  \draw[<->] (O)--node[left] {\(r_1\)} (S);
  \draw[<->] (O)--node[below right] {\(b\)} (B);

  \draw (E) ++(-0.52,-0.13) arc[start angle=194,end angle=227,radius=0.55];
  \node at (1.24,5.54) {\(\lvert\boldsymbol\epsilon\rvert\)};
  \draw (O) ++(0,0.75) arc[start angle=90,end angle=69,radius=0.75];
  \node at (0.26,0.92) {\(\vartheta\)};

  \fill (O) circle (1.3pt);
  \fill (S) circle (1.3pt);
  \node[below left] at (O) {\(r=0\)};
  \node[above left] at (S) {\(r=r_1\)};
\end{tikzpicture}
\caption{Quantities used in the calculation of parallaxes and apparent
luminosities. The angles and the curvature of the light ray are greatly
exaggerated.}
\label{fig:14.1}
\end{figure}
\endgroup


Here again, we use a vector notation, with
\[
  \mathbf{x}
  =
  \left(
    r\sin\theta\cos\phi,
    r\sin\theta\sin\phi,
    r\cos\theta
  \right),
\]
and we define scalar products as in Euclidean geometry. There is no loss of
generality in taking \(\mathbf{n}\) to be a unit vector, with
\(\mathbf{n}^2=1\). The parametric equation of the light paths, given by
substituting Eq.~\eqref{eq:14.4.1} in Eq.~\eqref{eq:14.4.2}, is then
\begin{equation}
  \mathbf{x}(\sigma)
  =
  \mathbf{n}\sigma
  +\mathbf{x}_1
  \left[
    \sqrt{1-k\sigma^2}
    -
    \left(1-\sqrt{1-kr_1^2}\right)
    \frac{(\mathbf{n}\cdot\mathbf{x}_1)\sigma}{r_1^2}
  \right],
  \tag{14.4.3}\label{eq:14.4.3}
\end{equation}
where \(r_1=(\mathbf{x}_1^2)^{1/2}\).

We shall now specify that the origin of the \(x^\mu\) coordinate system is
some definite point in the solar system, such as the center of the sun or the
center of the 200-in.\ mirror at Palomar, and we shall restrict ourselves to
light paths that pass close to this origin. In this case, the unit vector
\(\mathbf{n}\) must point nearly in the \(-\mathbf{x}_1\) direction, so
\begin{equation}
  \mathbf{n}\simeq-\widehat{\mathbf{x}}_1+\boldsymbol{\epsilon},
  \tag{14.4.4}\label{eq:14.4.4}
\end{equation}
where \(\widehat{\mathbf{x}}_1\equiv\mathbf{x}_1/r_1\), and
\(\boldsymbol{\epsilon}\) is a very small vector perpendicular to
\(\mathbf{x}_1\). (Here and below, \(\simeq\) means that an equation is
valid to first order in \(\boldsymbol{\epsilon}\).) Recalling
Eq.~\eqref{eq:14.4.1}, we note for future reference that
\(\lvert\boldsymbol{\epsilon}\rvert\) is the angle between the light path
and the \(-\mathbf{x}_1\) direction, as measured in the coordinate system
\(x'^\mu\), which is locally inertial at the light source. The light path,
given by inserting Eq.~\eqref{eq:14.4.4} in Eq.~\eqref{eq:14.4.3} and
discarding terms of order \(\boldsymbol{\epsilon}^2\), is
\begin{equation}
  \mathbf{x}(\sigma)
  \simeq
  -\widehat{\mathbf{x}}_1
  \left[
    \sigma\sqrt{1-kr_1^2}
    -r_1\sqrt{1-k\sigma^2}
  \right]
  +\boldsymbol{\epsilon}\sigma.
  \tag{14.4.5}\label{eq:14.4.5}
\end{equation}
The light path comes closest to the origin at \(\sigma\simeq r_1\). The
\emph{impact parameter} \(b\) is the proper distance of the path from the
origin at this point, given by Eqs.~\eqref{eq:14.2.1} and
\eqref{eq:14.4.5} as
\begin{equation}
  b
  \simeq
  a(t_0)\lvert\mathbf{x}(r_1)\rvert
  \simeq
  a(t_0)r_1\lvert\boldsymbol{\epsilon}\rvert,
  \tag{14.4.6}\label{eq:14.4.6}
\end{equation}
where \(t_0\) is the time that the light ray arrives near the origin.

Measurements of astronomical parallaxes amount to measurements of the
direction of light paths as a function of impact parameter, which in this
case is the projection of the earth--sun separation on the plane normal to
the line of sight. The light path has a direction near the origin given by
\[
  \left.
  \frac{\dd\mathbf{x}(\sigma)}{\dd\sigma}
  \right|_{\sigma\simeq r_1}
  \simeq
  \boldsymbol{\epsilon}
  -\frac{\widehat{\mathbf{x}}_1}{\sqrt{1-kr_1^2}},
\]
so the line of sight is given by a unit vector in the opposite direction
\begin{equation}
  \widehat{\mathbf{u}}
  \simeq
  -\sqrt{1-kr_1^2}
  \left.
  \frac{\dd\mathbf{x}(\sigma)}{\dd\sigma}
  \right|_{\sigma\simeq r_1}
  =
  \widehat{\mathbf{x}}_1
  -\sqrt{1-kr_1^2}\,\boldsymbol{\epsilon}.
  \tag{14.4.7}\label{eq:14.4.7}
\end{equation}
Hence the angle between the actual line of sight and the line of sight
\(\widehat{\mathbf{x}}_1\) that would be observed at the origin is
\begin{equation}
  \vartheta
  \simeq
  \lvert\widehat{\mathbf{u}}-\widehat{\mathbf{x}}_1\rvert
  \simeq
  \sqrt{1-kr_1^2}\,\lvert\boldsymbol{\epsilon}\rvert
  \simeq
  \sqrt{1-kr_1^2}\,
  \frac{b}{a(t_0)r_1}.
  \tag{14.4.8}\label{eq:14.4.8}
\end{equation}
In Euclidean geometry, a source at distance \(d\) would have a parallactic
angle \(\vartheta\simeq b/d\), so in general we may define the
\emph{parallax distance} \(d_P\) of a light source as
\begin{equation}
  d_P\equiv\frac{b}{\vartheta}
  \qquad
  \text{for \(\vartheta\to0,\ b\to0\),}
  \tag{14.4.9}\label{eq:14.4.9}
\end{equation}
and Eq.~\eqref{eq:14.4.8} may therefore be written
\begin{equation}
  d_P
  =
  a(t_0)\frac{r_1}{\sqrt{1-kr_1^2}}.
  \tag{14.4.10}\label{eq:14.4.10}
\end{equation}
In a universe with \(k=+1\), objects at \(r_1=1\) have infinite parallax
distance and further objects (with \(r_1<1\)) have decreasing parallax
distances, as noted first in 1900 by K.
Schwarzschild.\hyperref[ch14-ref-14]{[14]}

In order to calculate apparent luminosities, consider a circular telescope
mirror of radius \(b\), placed with its center at the origin and its normal
along the line of sight \(\widehat{\mathbf{x}}_1\) to the light source. The
light rays that just graze the mirror edge form a cone at the light source
that, in the coordinate system \(x'^\mu\) locally inertial at the source,
has a half-angle \(\lvert\boldsymbol{\epsilon}\rvert\) given by
Eq.~\eqref{eq:14.4.6}. The solid angle of this cone is
\[
  \pi\lvert\boldsymbol{\epsilon}\rvert^2
  =
  \frac{\pi b^2}{a^2(t_0)r_1^2},
\]
and the fraction of all isotropically emitted photons that reach the mirror
is the ratio of this solid angle to \(4\pi\), or
\begin{equation}
  \frac{\lvert\boldsymbol{\epsilon}\rvert^2}{4}
  =
  \frac{A}{4\pi a^2(t_0)r_1^2},
  \tag{14.4.11}\label{eq:14.4.11}
\end{equation}
where \(A\) is the proper area of the mirror,
\[
  A=\pi b^2.
\]
However, each photon emitted with energy \(h\nu_1\) will be red-shifted to
energy \(h\nu_1a(t_1)/a(t_0)\), and photons emitted at time intervals
\(\delta t_1\) will arrive at time intervals
\(\delta t_1a(t_0)/a(t_1)\), where as always \(t_1\) is the time the light
leaves the source, and \(t_0\) is the time the light arrives at the mirror.
Thus the total power \(P\) received by the mirror is the total power emitted
by the source, its absolute luminosity \(L\), times a factor
\(a^2(t_1)/a^2(t_0)\), times the fraction \eqref{eq:14.4.11}:
\[
  P
  =
  L
  \left[\frac{a^2(t_1)}{a^2(t_0)}\right]
  \left[\frac{A}{4\pi a^2(t_0)r_1^2}\right].
\]
The apparent luminosity \(l\) is the power per unit mirror area, so
\begin{equation}
  l
  \equiv
  \frac{P}{A}
  =
  \frac{L\,a^2(t_1)}{4\pi a^4(t_0)r_1^2}.
  \tag{14.4.12}\label{eq:14.4.12}
\end{equation}
In a Euclidean space the apparent luminosity of a source at rest at distance
\(d\) would be \(L/(4\pi d^2)\), so in general we may define the
\emph{luminosity distance} \(d_L\) of a light source as
\begin{equation}
  d_L\equiv\left(\frac{L}{4\pi l}\right)^{1/2},
  \tag{14.4.13}\label{eq:14.4.13}
\end{equation}
and Eq.~\eqref{eq:14.4.12} may therefore be written
\begin{equation}
  d_L=\frac{a^2(t_0)}{a(t_1)}\,r_1.
  \tag{14.4.14}\label{eq:14.4.14}
\end{equation}
(This calculation could also have been carried out without using quantum
theory, by applying the energy-conservation equation
\(\nabla_\nu T_\gamma^{\mu\nu}=0\) to the radiation emitted by the
source.\hyperref[ch14-ref-15]{[15]})

Next, let us calculate the angular diameter, observed at \(r=0,t=t_0\), of a
light source of true proper diameter \(D\) at \(r=r_1,t=t_1\). The light
rays from the edge of the source travel to the origin along fixed directions
\(\mathbf{x}_1/r_1\). Without loss of generality, we can rotate the
coordinate system so that the center of the light source is at \(\theta=0\),
and suppose that light from its edges travels to the origin on a cone with
half-angle \(\theta=\delta/2\). (See Figure~\ref{fig:14.2}.) The proper
distance across the source is then given by Eq.~\eqref{eq:14.2.1} as
\[
  D=a(t_1)r_1\delta
  \qquad\text{for \(\delta\ll1\),}
\]
so the angular diameter of the source is thus
\begin{equation}
  \delta=\frac{D}{a(t_1)r_1}.
  \tag{14.4.15}\label{eq:14.4.15}
\end{equation}

% Source: Weinberg GR, physical PDF pp. 442 and 444 (printed p. 422,
% repeated in the local scan).
% Coverage: Figure 14.2, angular diameter and proper-motion geometry.
% Status: vector reconstruction; source-reviewed and compile-clean.
% Geometry: the two edge rays reach the origin on a cone of half-angle
%           delta/2, as the accompanying text states, so the marked angle
%           delta spans both rays.  Here delta/2 = 8 degrees and the source
%           lies at radius 6.4, putting its edges at (-+0.8907, 6.3378).  The
%           proper diameter D is an arc length across the source, so it is
%           drawn as a circular arc of that radius about the origin, with the
%           straight measuring arrow set above it.
% Uncertainties: source caption states that the angle is deliberately
%                exaggerated.

\begingroup
\renewcommand{\thefigure}{14.2}
\begin{figure}[htbp]
\centering
\begin{tikzpicture}[x=1cm,y=1cm,>=Latex]
  \coordinate (O) at (0,0);
  \coordinate (L) at (-0.8907,6.3378);
  \coordinate (R) at (0.8907,6.3378);

  % Edge rays, travelling from the source down to the observer.
  \draw[line width=0.7pt,
    postaction={decorate},
    decoration={markings,mark=at position 0.45 with {\arrow{Latex}}}]
    (L) -- (O);
  \draw[line width=0.7pt,
    postaction={decorate},
    decoration={markings,mark=at position 0.45 with {\arrow{Latex}}}]
    (R) -- (O);

  % The source itself: proper diameter D, an arc length at radius r_1.
  \draw[line width=0.7pt] (98:6.4) arc[start angle=98,end angle=82,radius=6.4];
  \draw[line width=0.5pt] (0,6.4) -- (O);

  \draw[|<->|] (-0.8907,6.95) -- node[above,inner sep=2pt] {\(D\)}
               (0.8907,6.95);

  % The angular diameter subtended at the observer.
  \draw[<->] (-0.3232,2.30)
    -- node[above,inner sep=1.5pt,fill=white] {\(\delta\)} (0.3232,2.30);

  \node[right,inner sep=4pt] at (O) {\(r=0\)};
  \node[right,inner sep=4pt] at (R) {\(r=r_1\)};
\end{tikzpicture}
\caption{Quantities used in the calculation of angular diameters and proper
motions. The angle \(\delta\) is greatly exaggerated.}
\label{fig:14.2}
\end{figure}
\endgroup


In Euclidean geometry, the angular diameter of a source of diameter \(D\) at
a distance \(d\) is \(\delta=D/d\), so in general we may define the
\emph{angular diameter distance} \(d_A\) of a light source as
\begin{equation}
  d_A\equiv\frac{D}{\delta},
  \tag{14.4.16}\label{eq:14.4.16}
\end{equation}
and Eq.~\eqref{eq:14.4.15} may therefore be written
\begin{equation}
  d_A=a(t_1)r_1.
  \tag{14.4.17}\label{eq:14.4.17}
\end{equation}
Note that \(a(t_1)\) decreases as \(r_1\) increases, so in some models
\(d_A\) can have a maximum, with objects at very large distances having
angular diameters that increase with increasing luminosity distance.

Finally, let us consider the determination of distances from proper motions.
A source with a true velocity \(V_\perp\) transverse to the line of sight
will, in a time \(\Delta t_0\), move a proper distance
\[
  \Delta D
  =
  V_\perp\Delta t_1
  =
  V_\perp\Delta t_0\frac{a(t_1)}{a(t_0)},
\]
so, by the same reasoning that led to Eq.~\eqref{eq:14.4.15}, the source
will appear to move an angular distance
\begin{equation}
  \Delta\delta
  =
  \frac{\Delta D}{a(t_1)r_1}
  =
  \frac{V_\perp\Delta t_0}{a(t_0)r_1}.
  \tag{14.4.18}\label{eq:14.4.18}
\end{equation}
In a Euclidean space the change in apparent position on the celestial sphere
of a source at distance \(d\) would be \(V_\perp\Delta t_0/d\), so we may
define the \emph{proper-motion distance} of a light source as
\begin{equation}
  d_M\equiv\frac{V_\perp}{\mu},
  \tag{14.4.19}\label{eq:14.4.19}
\end{equation}
where \(\mu\) is the proper motion
\begin{equation}
  \mu\equiv\frac{\Delta\delta}{\Delta t_0}.
  \tag{14.4.20}\label{eq:14.4.20}
\end{equation}
Equation~\eqref{eq:14.4.18} may therefore be written
\begin{equation}
  d_M=a(t_0)r_1.
  \tag{14.4.21}\label{eq:14.4.21}
\end{equation}
Of course, we can use Eq.~\eqref{eq:14.4.19} to measure the proper-motion
distance only if we have some \emph{a priori} knowledge of the transverse
velocity, a point to which we return in the next section.

The luminosity distance \(d_L\), angular diameter distance \(d_A\), and
proper-motion distance \(d_M\) of a light source with red shift \(z\) are
related by the simple formulas
\begin{equation}
  \frac{d_A}{d_L}
  =
  \frac{a^2(t_1)}{a^2(t_0)}
  =
  (1+z)^{-2},
  \tag{14.4.22}\label{eq:14.4.22}
\end{equation}
\begin{equation}
  \frac{d_M}{d_L}
  =
  \frac{a(t_1)}{a(t_0)}
  =
  (1+z)^{-1}.
  \tag{14.4.23}\label{eq:14.4.23}
\end{equation}
If one can measure \(z\) accurately, there is no point in attempting a
separate determination of \(d_L,d_A,\) and \(d_M\), except perhaps as a
check of the Robertson--Walker metric or of the cosmological origin of red
shifts. In contrast, the measurement of the parallax distance \(d_P\) could
in principle give information beyond what could be learned from a
measurement of \(d_L\) and \(z\), but of course at present it is only
possible to measure parallaxes for very close objects, with \(z\ll1\) and
\(r_1\ll1\). In this case all these observable distances become essentially
equal to each other, and also to the proper distance \eqref{eq:14.2.21}:
\begin{equation}
  d_A
  \simeq d_L
  \simeq d_M
  \simeq d_P
  \simeq d_{\mathrm{prop}}(t_0)
  \simeq a(t_0)r_1.
  \tag{14.4.24}\label{eq:14.4.24}
\end{equation}
The distinction among different measures of distance only becomes important
for objects that are billions of light years away.

Actually, measurements of luminosity distance, angular diameter distance,
and red shift are inextricably mixed, for at least two reasons:

\medskip
\noindent\textit{(A) Light sources such as galaxies have smooth luminosity
distributions, without sharp edges.}
Let \(L(D)\) be the absolute luminosity of that part of a light source within
a circle (in the plane transverse to the line of sight) of diameter \(D\).
Then Eqs.~\eqref{eq:14.4.12} and \eqref{eq:14.4.15} give the apparent
luminosity within an angular diameter \(\delta\) as
\begin{equation}
  l(\delta)
  =
  \frac{
    L\!\left[r_1a(t_1)\delta\right]a^2(t_1)
  }{
    4\pi a^4(t_0)r_1^2
  }.
  \tag{14.4.25}\label{eq:14.4.25}
\end{equation}
It is more convenient to write this formula in terms of an absolute
luminosity per unit transverse area
\begin{equation}
  B(D)\equiv\frac{L'(D)}{2\pi D}
  \tag{14.4.26}\label{eq:14.4.26}
\end{equation}
and an apparent luminosity per unit solid angle, or \emph{brightness}:
\begin{equation}
  b(\delta)\equiv\frac{l'(\delta)}{2\pi\delta}.
  \tag{14.4.27}\label{eq:14.4.27}
\end{equation}
Using Eqs.~\eqref{eq:14.4.26}, \eqref{eq:14.4.27}, and
\eqref{eq:14.3.6} in Eq.~\eqref{eq:14.4.25} then gives the brightness as
\begin{equation}
  b(\delta)
  =
  \frac{
    B\!\left[r_1a(t_1)\delta\right]
  }{
    4\pi(1+z)^4
  }.
  \tag{14.4.28}\label{eq:14.4.28}
\end{equation}
The \emph{isophotal angular diameter} is the angle \(\delta_b\) at which the
brightness \eqref{eq:14.4.28} falls below some fixed threshold value \(b\):
\begin{equation}
  \delta_b=\frac{D_b}{r_1a(t_1)},
  \tag{14.4.29}\label{eq:14.4.29}
\end{equation}
where \(D_b\) is defined by the implicit equation
\begin{equation}
  B(D_b)=4\pi b(1+z)^4.
  \tag{14.4.30}\label{eq:14.4.30}
\end{equation}
For instance, Hubble has suggested that \(B(D)\) is well represented for
most galaxies by a function that near the galactic edge has the approximate
form\hyperref[ch14-ref-16]{[16]}
\begin{equation}
  B(D)\simeq\frac{\alpha L}{D^2},
  \tag{14.4.31}\label{eq:14.4.31}
\end{equation}
where \(\alpha\) is a dimensionless constant, of order unity. Then
Eqs.~\eqref{eq:14.4.29}--\eqref{eq:14.4.31} and
\eqref{eq:14.4.12} give
\begin{equation}
  D_b
  \simeq
  \left[
    \frac{\alpha L}{4\pi b(1+z)^4}
  \right]^{1/2}
  \tag{14.4.32}\label{eq:14.4.32}
\end{equation}
and
\begin{equation}
  \delta_b
  \simeq
  \left(\frac{\alpha l}{b}\right)^{1/2}.
  \tag{14.4.33}\label{eq:14.4.33}
\end{equation}
In this particular case, the measurement of an isophotal angular diameter is
just tantamount to a measurement of apparent luminosity.

\medskip
\noindent\textit{(B) Most detectors of radiation respond only to photons in
a narrow range of wavelengths.}
Thus it is necessary to distinguish between the \emph{bolometric
luminosities} \(L\) or \(l\) discussed above, which take account of
radiation emitted or received at all wavelengths, and the
\emph{ultraviolet, blue, photographic, visual, and infrared luminosities},
which represent the average power or flux in various wavelength bands. If a
source emits a radiant power \(L(\nu_1)\) at all frequencies less than
\(\nu_1\), then Eqs.~\eqref{eq:14.4.12} and \eqref{eq:14.3.4} give an
apparent luminosity for all frequencies less than \(\nu_0\) as
\begin{equation}
  l(\nu_0)
  =
  \frac{
    L\!\left[\nu_0a(t_0)/a(t_1)\right]a^2(t_1)
  }{
    4\pi a^4(t_0)r_1^2
  }.
  \tag{14.4.34}\label{eq:14.4.34}
\end{equation}
The frequency distributions of received and emitted power are therefore
related by the formula
\begin{equation}
  l'(\nu_0)
  =
  \frac{
    L'\!\left[\nu_0a(t_0)/a(t_1)\right]a(t_1)
  }{
    4\pi a^3(t_0)r_1^2
  }.
  \tag{14.4.35}\label{eq:14.4.35}
\end{equation}
For a black body, \(L'(\nu)\) is given by the Planck formula
\begin{equation}
  L'(\nu)
  =
  \frac{15L}{\pi^4\nu}
  \left(\frac{h\nu}{k_{\mathrm B}T_1}\right)^4
  \left[
    \exp\!\left(\frac{h\nu}{k_{\mathrm B}T_1}\right)-1
  \right]^{-1},
  \tag{14.4.36}\label{eq:14.4.36}
\end{equation}
where \(T_1\) is the source temperature, \(k_{\mathrm B}\) is the Boltzmann
constant, and \(h\) is Planck's constant. The frequency distribution of
received radiation is then
\begin{equation}
  l'(\nu_0)
  =
  \frac{15l}{\pi^4\nu_0}
  \left(\frac{h\nu_0}{k_{\mathrm B}T_0}\right)^4
  \left[
    \exp\!\left(\frac{h\nu_0}{k_{\mathrm B}T_0}\right)-1
  \right]^{-1},
  \tag{14.4.37}\label{eq:14.4.37}
\end{equation}
where \(l\) is the bolometric apparent luminosity \eqref{eq:14.4.12}, and
\(T_0\) is the red-shifted temperature
\begin{equation}
  T_0=T_1\frac{a(t_1)}{a(t_0)}.
  \tag{14.4.38}\label{eq:14.4.38}
\end{equation}
If we know the temperatures \(T_1\) or \(T_0\), it is easy to convert the
absolute luminosity \(L'(\nu_1)\Delta\nu_1\) or the apparent luminosity
\(l'(\nu_0)\Delta\nu_0\) in any narrow frequency band into a bolometric
absolute or apparent luminosity.

A word must be said about the time-honored language used by astronomers to
describe astronomical distances and luminosities. The astronomical unit
(abbreviated a.u.) is the mean distance of the earth from the sun:
\begin{equation}
  1\ \mathrm{a.u.}
  =
  1.49598\times10^8\ \mathrm{km}.
  \tag{14.4.39}\label{eq:14.4.39}
\end{equation}
Treating the earth's orbit as a circle, the projection of the earth--sun
separation vector on the plane normal to the line of sight to any fixed star
reaches a maximum value \(b_{\max}\) equal to \(1\ \mathrm{a.u.}\) during
the course of a year, so the position of a star traces out an ellipse of
maximum radius \(\varpi\) given by Eq.~\eqref{eq:14.4.9} as
\begin{equation}
  \varpi\ \text{(in radians)}
  =
  \frac{1}{d_P\ \text{(in a.u.)}}.
  \tag{14.4.40}\label{eq:14.4.40}
\end{equation}
We shall call \(\varpi\) the \emph{trigonometric parallax}. One parsec
(abbreviated pc) is defined as the distance \(d_P\) at which a star would
have a trigonometric parallax of \(1''\); there are \(206{,}264.8\) seconds
in one radian, so
\begin{equation}
  \begin{aligned}
  1\ \mathrm{pc}
  &=
  206{,}264.8\ \mathrm{a.u.}
  =
  3.0856\times10^{13}\ \mathrm{km}
  \\
  &=
  3.2615\ \text{light years}.
  \end{aligned}
  \tag{14.4.41}\label{eq:14.4.41}
\end{equation}
Thus Eq.~\eqref{eq:14.4.40} may in general be expressed as
\begin{equation}
  \varpi\ \text{(in seconds)}
  =
  \frac{1}{d_P\ \text{(in pc)}}.
  \tag{14.4.42}\label{eq:14.4.42}
\end{equation}
Only the nearest stars have measurable trigonometric parallaxes, but such is
the power of tradition that all astronomical distances outside our solar
system are conventionally given in parsecs, and sometimes these distances,
however measured, are even described in terms of an equivalent parallax.

The apparent bolometric luminosity \(l\) is usually expressed in terms of an
apparent bolometric magnitude \(m_{\mathrm{bol}}\), or simply \(m\), which
for historical reasons is defined so that
\begin{equation}
  l
  =
  10^{-2m/5}
  \times
  2.52\times10^{-5}\ 
  \mathrm{erg\,cm^{-2}\,s^{-1}}.
  \tag{14.4.43}\label{eq:14.4.43}
\end{equation}
The absolute bolometric magnitude \(M\) is defined as the apparent
bolometric magnitude the source would have at a distance \(10\ \mathrm{pc}\),
so
\begin{equation}
  L
  =
  10^{-2M/5}
  \times
  3.02\times10^{35}\ 
  \mathrm{erg\,s^{-1}}.
  \tag{14.4.44}\label{eq:14.4.44}
\end{equation}
Equation~\eqref{eq:14.4.13} may be expressed as a formula for the luminosity
distance \(d_L\) in terms of the \emph{distance modulus} \(m-M\):
\begin{equation}
  d_L=10^{1+(m-M)/5}\ \mathrm{pc}.
  \tag{14.4.45}\label{eq:14.4.45}
\end{equation}
The apparent magnitudes \(m_U,m_B,\) and so on, in the ultraviolet, blue,
photographic, visual, and infrared wavelength bands are related to the
corresponding apparent luminosities by formulas like
Eq.~\eqref{eq:14.4.43}, but with different normalization constants, chosen
so that all apparent magnitudes will be the same for stars of the A0 spectral
type between fifth and sixth magnitude. The corresponding absolute
magnitudes are defined so that the distance moduli \(m_U-M_U\),
\(m_B-M_B\), and so on, are all equal to \(m-M\). (Often the ultraviolet,
blue, and visual apparent magnitudes \(m_U,m_B,m_V\) are denoted \(U,B,V\).)
The \emph{color index} is the quantity
\(m_B-m_V=M_B-M_V\); stars with negative color index are bluer than stars
with positive color index. For purposes of comparison, the sun has absolute
magnitudes
\[
  M(\text{bolometric})=+4.72,
  \qquad
  M_U=5.51,
  \qquad
  M_B=5.41,
  \qquad
  M_V=4.79,
\]
and apparent magnitudes
\[
  m(\text{bolometric})=-26.85,
  \qquad
  m_U=-26.06,
  \qquad
  m_B=-26.16,
  \qquad
  m_V=-26.78,
\]
so that its distance modulus is \(-31.57\) and its color index is \(0.62\).
